\hyperref[ch:032]{``群与对称：可逆操作怎样组织不变性''}把对称操作组织成了群，交出来的成果是一张乘法表。表格是精确的，也是不能计算的：格子里装着操作的名字，名字不能相加，没有特征值，也谈不上远近。想比较两个群、想问一个群内部有没有几套彼此不干扰的机制，从乘法表出发都无从下手。

本章的做法是给每个群元素配一个可逆矩阵，让群的复合对应矩阵乘法。这一步之后，对称性住进了向量空间，线性代数的全套工具——迹、特征值、不变子空间、直和分解——立刻可以用来研究它。

真正的收获是分解。Maschke 定理说，在合适的系数域上（有限群、域的特征不整除群阶），任何表示都能拆成不可约表示的直和；换一组基，整族矩阵同时变成分块对角，块与块之间没有信息流动。不可约表示因此是对称性的原子。这与谱分解是同一个念头的不同版本：把一个复杂对象拆成互不干扰的基本模式，只不过谱定理处理单个算子，这里处理一整族。

配套的是不变量。特征标定义为表示矩阵的迹，它在相似变换下不变，因而不依赖基的选择；不可约特征标彼此正交，于是``这个表示含有哪些成分、各多少份''被化简成一次内积计算，不必再去矩阵里搜索不变子空间。分解与不变量这两件事，构成本章的全部主线。

对做数据的人，这条线有几个直接的出口。等变神经网络要求层映射与群作用交换，Schur 引理把这类映射的形式限制得极死：不同不可约成分之间只能是零，相同成分之间只剩一个标量，参数共享因此是被代数逼出来的而非人工设计；三维点云与分子建模里常用的球谐函数，正是 \(SO(3)\) 的不可约表示的基。交换群上的表示论则直接就是 Fourier 分析——离散 Fourier 变换的那组指数基，是循环群仅有的不可约表示。分块对角在计算上意味着大问题裂成互不干扰的小问题，卷积定理是它最常见的一副面孔。

\subsection{本章篇目}

\begin{itemize}
\tightlist
\item \hyperref[sec:04-Discrete-Mathematics/09-Representation-Theory/01]{``表示把群变成矩阵群''}：从正三角形的六个动作出发，说清同态这一条要求带来了什么，以及等价表示为何只是统一换基；
\item \hyperref[sec:04-Discrete-Mathematics/09-Representation-Theory/02]{``不可约表示是基本构建块''}：不变补空间未必存在，Maschke 定理靠平均化把它造出来，条件塌掉时反例立刻回来；
\item \hyperref[sec:04-Discrete-Mathematics/09-Representation-Theory/03]{``特征标是表示的指纹''}：迹丢掉了基坐标却保住了分解信息，正交关系把分解变成求坐标。
\end{itemize}

读这一章需要群、共轭类与群作用的基本概念（见\hyperref[sec:04-Discrete-Mathematics/06-Groups-and-Symmetry/02]{``子群、陪集与群作用''}），以及相似变换与对角化。第一遍读完能回答三件事就够了：群怎样作用在向量空间上，哪些子空间被所有群元素共同保持，以及手上的结论依赖有限群、复数域还是完全可约性中的哪一条。
